\chapter{Trang quản trị và cổng API}

\section{Trang quản trị React}

\subsection{Công nghệ}

\begin{longtable}{L{4.4cm} L{9.2cm}}
\toprule
\textbf{Thành phần} & \textbf{Lựa chọn}\\
\midrule
\endhead
Khung & React 19 + TypeScript, dựng bằng Vite 8\\
Thư viện giao diện & Ant Design 6\\
Truy vấn dữ liệu & TanStack Query 5 — quản lý bộ đệm, thử lại, làm mới nền\\
Biểu mẫu & React Hook Form + bộ giải kiểm tra hợp lệ\\
Hoạt ảnh & GSAP 3 kèm \code{@gsap/react}\\
HTTP & Axios\\
Xác thực & JWT chung với backend, cộng đăng nhập quản trị có mã một lần (OTP)\\
Triển khai & Vercel\\
\bottomrule
\caption{Công nghệ trang quản trị}
\end{longtable}

\subsection{Hai mươi chín trang chức năng}

\begin{longtable}{L{4.0cm} L{9.6cm}}
\toprule
\textbf{Nhóm} & \textbf{Trang}\\
\midrule
\endhead
Bảng điều khiển & \code{EnhancedAdminDashboard}, \code{SuperAdminDashboard}, \code{RoleRedirectPage} (điều hướng theo vai trò).\\
Nội dung học & \code{CoursesPage}, \code{UnitsPage}, \code{LessonsPage}, \code{LessonExercisesPage}, \code{VocabularyPage}, \code{TopicsPage}.\\
Tác tử nội dung & \code{ContentLabPage} (chạy tác tử sinh khoá học), \code{ContentQaQueuePage} (hàng đợi kiểm duyệt), \code{ContentAnalyticsPage}.\\
Người dùng & \code{UserManagementPage}, \code{AdminManagementPage}, \code{NoAccessPage}.\\
Động lực học & \code{AchievementsPage}, \code{ShopPage}, \code{RankingAgentPage}.\\
AI & \code{AiModelsPage}, \code{AiChatSettingsPage}, \code{AiQualityDashboardPage}.\\
Vận hành & \code{MonitoringPage}, \code{DatabasePage}, \code{LogsPage}, \code{SystemSettingsPage}.\\
Tăng trưởng & \code{NotificationCampaignPage}, \code{AdsPage}.\\
Khác & \code{LoginPage}, \code{NotFoundPage}.\\
\bottomrule
\caption{Trang của bảng điều khiển quản trị}
\end{longtable}

\subsection{Quy trình kiểm duyệt nội dung}

Đây là quy trình quan trọng nhất của trang quản trị, vì nó là điểm con người can thiệp vào
nội dung do AI sinh:

\begin{enumerate}[leftmargin=1.4em]
  \item Quản trị viên tải lên danh sách từ có nhãn CEFR ở \code{ContentLabPage}.
  \item Tác tử nội dung chạy, kết quả lưu thành khoá học \textbf{bản nháp}.
  \item \code{ContentQaQueuePage} liệt kê bài học cần rà soát; người duyệt đọc nội dung,
  sửa hoặc yêu cầu sinh lại.
  \item Chỉ khi mọi bài học có bài tập, nút xuất bản mới mở
  (\code{CourseCRUD.publish\_blockers}).
\end{enumerate}

\section{Cổng API}

\subsection{Nginx}

Cổng sản xuất là Nginx 1.27-alpine, cấu hình theo khuôn mẫu trong
\code{gateway/nginx/templates/} kèm các đoạn dùng lại ở \code{snippets/}. Trách nhiệm:

\begin{itemize}[leftmargin=1.4em]
  \item Kết thúc TLS và chuyển hướng HTTP sang HTTPS.
  \item Định tuyến \code{/api/v1/*} về backend-service, \code{/ai/*} và các đường hội thoại
  về ai-service.
  \item Nâng cấp giao thức cho WebSocket (giọng nói thời gian thực).
  \item Giới hạn tần suất ở tầng biên, trước khi yêu cầu chạm ứng dụng.
  \item Chuyển tiếp \code{X-Request-ID} và địa chỉ IP thật.
\end{itemize}

Kho mã còn giữ cấu hình Kong (\code{gateway/kong/},
\code{docker-compose.kong.yml}) như phương án thay thế và cấu hình APIM cho môi trường
lai.

\subsection{Tường lửa ứng dụng web ở Cloudflare}

Bộ quy tắc WAF (\code{gateway/cloudflare/waf-rules.md}) được rút ra từ phân tích nhật ký
30 ngày. Quy tắc quan trọng nhất phải đặt \textbf{đầu tiên}:

\begin{lstlisting}
(http.request.method eq "OPTIONS") and (http.host eq "api.lexilingo.me")
  -> Skip: Bot Fight Mode
\end{lstlisting}

\begin{luuy}[Vì sao quy tắc này phải đứng đầu]
Chế độ chống bot của Cloudflare chặn yêu cầu \code{OPTIONS} (yêu cầu tiền kiểm CORS)
\textbf{trước khi} nó tới backend và trả về \code{cf-mitigated: challenge}. Hậu quả là mọi
lời gọi API có xác thực từ \code{www.lexilingo.me} sang \code{api.lexilingo.me} đều hỏng —
triệu chứng nhìn giống lỗi CORS ở phía ứng dụng, nên rất dễ đi tìm sai chỗ.
\end{luuy}

Các nhóm quy tắc còn lại: chặn quét lỗ hổng theo đường dẫn quen thuộc, giới hạn tần suất
theo IP cho điểm cuối đăng nhập, thách thức lưu lượng bất thường, và danh sách cho phép
đối với các dịch vụ nội bộ.

\section{Máy chủ MCP}

\code{mcp-server/} triển khai \en{Model Context Protocol} — giao thức cho phép trợ lý lập
trình trong IDE gọi công cụ của dự án. Các công cụ hiện có:

\begin{longtable}{L{4.4cm} L{9.2cm}}
\toprule
\textbf{Công cụ} & \textbf{Chức năng}\\
\midrule
\endhead
\code{manage\_i18n\_key} & Thêm hoặc cập nhật một khoá dịch đồng thời trên mọi tệp ngôn ngữ.\\
\code{whisper\_handler} & Chạy nhận dạng giọng nói cục bộ.\\
\code{hubert\_handler} & Phân tích âm vị cục bộ.\\
\code{piper\_handler} & Tổng hợp giọng nói cục bộ.\\
\code{qwen\_handler}, \code{ollama\_qwen\_handler} & Suy luận Qwen cục bộ.\\
\code{gemini\_handler} & Gọi API Gemini.\\
\bottomrule
\caption{Công cụ của máy chủ MCP}
\end{longtable}

Đây là thành phần phục vụ đội phát triển chứ không nằm trên đường phục vụ người học, nên
nó không có mặt trong sơ đồ triển khai sản xuất.
